\bookchapter{第六十五章：项目一——我的记账本（本地版）}

\chapterintro{欢迎来到你的第一个“正经项目”。这个记账本虽然不联网，但五脏俱全：能记一笔账、能看流水、能按月统计、能删改。做完它，你就完整走了一遍“数据 逻辑 界面”的全流程。别急着抄代码——\textbf{先想清楚再动手}，这是咱们一贯的原则。}

\booksection{一、需求思考：先问五个问题}

\begin{center}
\begin{tabularx}{\linewidth}{|X|X|}
\hline
\textbf{要素} & \textbf{答案} \tabularnewline
\hline
Who（给谁用） & 自己用，家人朋友也行 \tabularnewline
What（做什么） & 记收支、看流水、按月汇总、按分类统计 \tabularnewline
Why（为什么） & 练手 + 搞清楚钱花哪儿了 \tabularnewline
How（怎么做） & Swift + SwiftUI + SwiftData + MVVM，纯本地无后端 \tabularnewline
Constraint（约束） & 一周内完成；最低支持 iOS 17（SwiftData 门槛） \tabularnewline
\hline
\end{tabularx}
\end{center}

\begin{center}\small\color{gray}记账本需求五要素\end{center}

然后拆解成“每一步都能用”的小步骤：(1) 建 SwiftData 数据层（能存一笔账）；(2) 列表页（能看流水）；(3) 录入弹窗（能加账）；(4) 删除与编辑；(5) 月度汇总卡片。每完成一步，App 都是一个“能用的东西”。

\booksection{二、数据层：让“账”有个家}

先定义一笔账长什么样——字段要够用但别贪多：

\begin{lstlisting}[language=Swift]
import SwiftData

@Model
class Record {
    var type: String      // "income" 收入 | "expense" 支出
    var amount: Double    // 金额（元）
    var category: String  // 分类：餐饮、交通、工资……
    var note: String      // 备注（可选）
    var date: Date        // 记账日期

    init(type: String, amount: Double, category: String, note: String = "", date: Date = Date()) {
        self.type = type
        self.amount = amount
        self.category = category
        self.note = note
        self.date = date
    }
}
\end{lstlisting}

\begin{tcbinfo}{小贴士：金额为什么用 Double？}
严谨的金融 App 会用 Decimal（精确十进制）来避免浮点误差（0.1+0.2$\neq$0.3 的老问题）。练手项目用 Double 足够；如果哪天做真金融产品，记得换成\texttt{Decimal}或“以分为单位的 Int”——这是无数前辈踩出来的坑。

\end{tcbinfo}

\booksection{三、MVVM：界面与逻辑“分家”}

直接让界面操作数据库？能跑，但界面代码会又臭又长。MVVM 架构让各司其职：\textbf{View（界面）只管画，ViewModel（逻辑中枢）管数据，Model（SwiftData）管存储}。

View（SwiftUI 界面）

ViewModel（@Observable 逻辑中枢）

Model（SwiftData 数据层）

\begin{center}\small\color{gray}图 4：MVVM 三层结构——数据单向流动，各层互不越界\end{center}

\begin{lstlisting}[language=Swift]
import Observation

@Observable
class RecordViewModel {
    var records: [Record] = []
    private let context: ModelContext

    init(context: ModelContext) { self.context = context }

    func load() {
        // SwiftData 的 #Predicate 筛选 + 排序
        let descriptor = FetchDescriptor<Record>(
            sortBy: [SortDescriptor(\.date, order: .reverse)]
        )
        records = (try? context.fetch(descriptor)) ?? []
    }

    func add(type: String, amount: Double, category: String, note: String) {
        let r = Record(type: type, amount: amount, category: category, note: note)
        context.insert(r)
        try? context.save()
        load()                       // 重新加载，界面刷新
    }

    func delete(_ record: Record) {
        context.delete(record)
        try? context.save()
        load()
    }
}
\end{lstlisting}

界面里\texttt{@State private var vm = RecordViewModel(.\allowbreak .\allowbreak .\allowbreak )}，然后直接\texttt{vm.\allowbreak records}渲染列表——数据一变，界面自动刷新。

\booksection{四、界面：列表页 + 录入弹窗}

\begin{lstlisting}[language=Swift]
struct RecordListView: View {
    @State private var vm: RecordViewModel
    @State private var showAdd = false

    var body: some View {
        NavigationStack {
            List {
                Section("本月结余") {          // 汇总卡片
                    SummaryRow(records: vm.records)
                }
                Section("流水") {
                    ForEach(vm.records) { r in
                        RecordRow(record: r)
                    }
                    .onDelete { idx in
                        idx.forEach { vm.delete(vm.records[$0]) }
                    }
                }
            }
            .navigationTitle("我的记账本")
            .toolbar {
                Button { showAdd = true } label: {
                    Image(systemName: "plus")
                }
            }
            .sheet(isPresented: $showAdd) {
                AddRecordView(vm: vm)          // 录入弹窗
            }
        }
    }
}
\end{lstlisting}

录入弹窗要点：金额输入框键盘用\texttt{.\allowbreak keyboardType(.\allowbreak decimalPad)}（数字键盘），收入/支出用分段控件（Picker），分类用菜单选择，点“保存”调\texttt{vm.\allowbreak add(.\allowbreak .\allowbreak .\allowbreak )}并\texttt{dismiss()}关窗。

\begin{tcbwarn}{注意：输入校验别偷懒}
金额输入框一定要校验：空值、负数、“abc”这类非法输入都要拦住（用\texttt{Double(text)}判断，失败返回 nil）。App 崩掉最常见的原因之一就是“把用户的输入当成了乖孩子”。记住：\textbf{用户输入的一切，都要当敌人审查}。

\end{tcbwarn}

\booksection{五、编辑、删除与汇总}

\begin{itemize}
  \item \textbf{删除}：List 的 .onDelete 左滑删除 + 确认弹窗（误删是记账类 App 用户最痛恨的事）；
  \item \textbf{编辑}：点行进编辑弹窗，预填原值（改一个字比重新输入省事）；
  \item \textbf{汇总}：用\texttt{filter}+\texttt{reduce}算"本月收入 / 支出 / 结余"，顶部卡片展示：
\end{itemize}

\begin{lstlisting}[language=Swift]
func monthSummary(_ records: [Record]) -> (income: Double, expense: Double) {
    let cal = Calendar.current
    let monthRecords = records.filter {
        cal.isDate($0.date, equalTo: Date(), toGranularity: .month)
    }
    let income = monthRecords
        .filter { $0.type == "income" }
        .reduce(0) { $0 + $1.amount }
    let expense = monthRecords
        .filter { $0.type == "expense" }
        .reduce(0) { $0 + $1.amount }
    return (income, expense)
}
\end{lstlisting}

\booksection{六、验收：三层检查 + 两个进阶挑战}

项目做完，按“功能 边界 架构”三层验收：能不能记、能不能删、重启 App 数据还在不在（功能）；金额输 0 或负数会不会崩（边界）；界面和数据有没有混在一起、字段有没有硬编码（架构）。

基础版验收通过后，给自己加两道“进阶菜”：

\begin{enumerate}
  \item \textbf{分类统计}：加一个"统计"标签页，用\textbf{Swift Charts}把分类支出画成饼图/柱状图——苹果官方图表库，几行代码出图；
  \item \textbf{iCloud 同步}：SwiftData 开 iCloud 同步只需配几项设置——账本在 iPhone/iPad 之间自动同步，瞬间"高级"起来。
\end{enumerate}

\begin{tcbpractice}{来动手：本周完成它}
按“数据层 ViewModel 列表页 录入弹窗 汇总卡片”的顺序动手，每天一步。卡住了就把报错原样丢给搜索引擎或 AI（还记得第七篇《AI编程实战篇》吗？描述问题时说清“期望行为、实际行为、复现步骤”）。完成后用它记一周账——自己做的 App 用起来，成就感会推着你继续往前走。

\end{tcbpractice}

\begin{tcbsummary}{本章小结}
\begin{itemize}
  \item 需求五要素 拆成"每步能用"的小步骤，再动手
  \item 数据层：@Model 模型 + modelContainer 容器 + FetchDescriptor 查询
  \item MVVM 分工：View 只管画，@Observable 的 ViewModel 管逻辑，SwiftData 管存储
  \item 界面：List + Section + toolbar + sheet 录入弹窗
  \item 三层验收（功能/边界/架构）+ 进阶挑战（Swift Charts、iCloud 同步）
\end{itemize}

\end{tcbsummary}
